\section{Results}
\subsection{The serial control isolates the effect of component selection}
In P0's primary routing-feedback setting, retained-load planning reduces mean cost by 38.2\% relative to static planning and by 38.1\% relative to capacity-at-20 (Table~\ref{tab:primary}). The paired, multiplicity-adjusted 97.5\% intervals are $[-51.56,-19.40]$ and $[-51.23,-19.32]$ thousand. Relative to predictive selection, the reduction is 35.4\%, with descriptive 95\% interval $[-46.30,-16.66]$ thousand.

\begin{table}[t]\caption{P0: mean realized cost in thousands over 12 independent fits per mechanism. All methods use the same test conditions.}\label{tab:primary}\centering\small\setlength{\tabcolsep}{3pt}
\begin{tabular}{lrrrrrr}\toprule
Mechanism & \shortstack{No\\action} & Static & Retained & Predictive & \shortstack{Capacity\\at 20} & \shortstack{Rework\\at 20}\\\midrule
Fixed & 23.6 & 18.0 & 19.9 & 18.0 & 17.7 & 22.6 \\
Service & 211.7 & 111.0 & 111.0 & 111.0 & 111.0 & 202.7 \\
Routing & 147.8 & 92.8 & 57.4 & 88.8 & 92.6 & 45.3 \\
Combined & 375.2 & 309.2 & 288.3 & 292.5 & 309.2 & 284.1 \\
\bottomrule

\end{tabular}\end{table}

Predictive selection retains the return-load component in 2 of 12 fits and chooses rework reduction in 1; retained-load planning chooses it in 9. Every retained-load fit selects threshold 20, so this result concerns action choice, not improved timing. The fixed rework rule and known-generator planner both cost 45.3 thousand, below retained-load's 57.4. Their paired difference is $12.04$ thousand, with 95\% interval $[-1.87,25.94]$. A feedback-strength heuristic costs 73.0 thousand and chooses rework in 5 fits; the retained-load minus heuristic interval is $[-38.90,7.65]$ thousand. These comparisons do not establish an advantage over either strong simple baseline.

The known-generator planner selects rework-at-20 for all 12 fits. Retained-load's excess cost comes from its 3 capacity choices: their fitted models reverse the ranking observed on the independent test. Those calibration histories reach only 12--19 unfinished cases. The error therefore concerns action ranking beyond observed support; the experiment does not isolate a unique incorrect coefficient.

The other mechanisms distinguish useful adaptation from indiscriminate complexity. Under service feedback all planners choose capacity-at-20. Under fixed behavior, retained-load raises cost by 10.9\% relative to static planning; under combined feedback it reduces cost by 6.7\%, while severe-event incidence increases from 69.3\% to 71.0\%. A lower integrated cost therefore need not imply fewer sustained overload events.

\subsection{Selector sensitivity and observable response behavior}
The P0 diagnostic reuses every saved test action while varying the return-selection threshold over $0,0.001,0.0025,0.005,0.01$. The selector retains the component in $8,5,2,2,0$ of the 12 routing fits, respectively; mean costs are $80.8,84.9,88.8,88.8,92.8$ thousand. Retained-load remains at $57.4$ thousand. The observed difference is consequently not confined to the original threshold $0.005$. Varying the duration threshold over $0,0.0025,0.005,0.01,0.02$ gives selector costs of $85.1$ thousand at zero and $88.8$ at the other settings. These are diagnostics on the existing independent-test catalog, not additional confirmatory datasets. Duration-only and return-only retention cost $69.4$ and $80.8$ thousand, respectively; both retained components produce the lowest mean among these ablations.

In all 289 P0 routing trajectories with a severe event under no action, the retained-load rule activates beforehand, with mean lead 108.0. It prevents 93 events (32.2\%). No activation occurs in 192 calm trajectories, but 95 of the 287 trajectories without a no-action severe event activate the rule. With the originally selected policies unchanged, actual rework reductions of 30\%, 50\% and 70\% yield cost reductions of 17.0--48.7\% relative to static planning. These results separate timely activation, event prevention, cost reduction and robustness to action-effect error.

\subsection{How structure changes the value of retaining load dependencies}
All 288 new fits complete, and all 3168 calibration/validation episodes satisfy the pre-crisis criterion. Table~\ref{tab:structures_results} compares routing-feedback decisions with supplied structure. Retained-load reduces mean cost relative to predictive selection by 8.1\% in P1, 29.8\% in P2, 34.6\% in P5 and 26.4\% in P6. Cost increases by 3.0\% in P3 and 15.7\% in P4. None of the four improvements excludes zero after correction for six primary contrasts; P4's increase does, with an adjusted difference interval of $[1.9,8.3]$ thousand cost units. The extension therefore does not establish a per-structure confirmatory improvement.

\begin{table}[t]\caption{Routing feedback, known structure: mean per-fit cost divided by matched no-action cost. Regret is retained-load minus known-generator cost, normalized by no action. Each row contains 12 independent fits.}\label{tab:structures_results}\centering\small\setlength{\tabcolsep}{4pt}
\begin{tabular}{lrrrrrr}\toprule
Process & Static & Predictive & Retained & Heuristic & Generator & Regret\\\midrule
P1 & 0.587 & 0.549 & 0.502 & 0.501 & 0.429 & 0.073 \\
P2 & 0.731 & 0.614 & 0.428 & 0.514 & 0.428 & 0.000 \\
P3 & 0.630 & 0.630 & 0.650 & 0.630 & 0.516 & 0.133 \\
P4 & 0.520 & 0.528 & 0.610 & 0.620 & 0.512 & 0.098 \\
P5 & 1.030 & 0.727 & 0.473 & 0.422 & 0.422 & 0.050 \\
P6 & 0.575 & 0.575 & 0.422 & 0.460 & 0.349 & 0.073 \\
\bottomrule\end{tabular}\end{table}


The secondary equal-weight comparison shows where the average benefit lies. Under routing feedback, retained-minus-predictive normalized cost is $-0.090$ with descriptive 95\% interval $[-0.135,-0.045]$; under combined feedback it is $-0.036$ $[-0.052,-0.020]$. Fixed and service-only mechanisms instead increase by 0.058 and 0.069. These are differences in cost divided by each fit's no-action cost, not percentage reductions relative to the predictive planner. The result supports conditional value for retaining load dependencies in this designed family, rather than a rule to retain every load component.

\begin{figure}[t]\centering
\includegraphics[width=\textwidth]{generated/structural_matrix_en.pdf}
\caption{Mean cost reduction of retained-load relative to predictive selection, in percent. Positive values favor retention. Panels separate supplied and discovered structure. The asterisk marks a nonzero primary contrast after the six-comparison correction; other cells are descriptive.}\label{fig:structural_matrix}
\end{figure}

\subsection{Discovery errors can change the prepared action}
With discovered structure, the routing aggregate becomes $0.001$ $[-0.037,0.038]$: the known-structure advantage does not carry through the complete pipeline. P2 illustrates the decision consequence. With supplied structure, retained-load chooses the generator planner's rework action in all 12 fits and achieves normalized cost 0.428. With discovered structure, it chooses capacity in every fit and costs 0.897. The training data, learned behavioral parameters, resources and action catalog are unchanged between these variants. Thus, the structural representation changes which intervention the same behavioral estimates support. P1, P2 and P6 add paths in every discovered model despite fitness 1.0; mean precision is 0.497, 0.642 and 0.487. P3, P4 and P5 recover equivalent visible languages and reproduce the supplied-structure policies in all 48 fits per process. Trace fit alone therefore does not establish decision fidelity.

\subsection{What simulation adds beyond a local feedback rule}
P3 and P5 provide three distinct intervention types, but their comparisons do not establish a retained-load advantage over the feedback-strength rule. In P5 routing feedback, the rule always selects rework-at-10 and essentially matches generator planning; retained-load selects that action type in 11 of 12 fits. In P3, generator planning always targets the local return, whereas retained-load does so in only 1 of 12 fits despite exact control flow. The gap is in the learned action ranking. In P4, removing learned routing load terms reduces normalized cost from 0.610 to 0.516; removing duration load terms leaves it at 0.611. This ablation locates the error in the estimated return dependencies, without selecting a revised method from test outcomes.

The P5 replay makes the resource interaction explicit. Under combined feedback, adding a B worker lowers its mean queue from 141.02 to 40.19, but increases D's queue from 1.49 to 118.56. The dominant queue moves from B to D in all 24 matched diagnostic paths. Accelerating the upstream stage feeds work into a downstream stage whose service and repetition behavior also depend on load. These queue measurements use the original test horizon and fixed action rules.

P6 service feedback gives a different diagnostic result. The local rule usually adds capacity to the parallel branch; simulation-based retained-load planning more often chooses the shared pool and has normalized cost 0.633 versus 0.920. The descriptive paired cost interval is $[-26.3,-7.0]$ thousand. Static planning has a lower mean of 0.494, close to the generator's 0.492. This comparison supports assessing system-wide consequences when local feedback scores mislead, while showing that the additional load components are unnecessary for that particular improvement.

\subsection{Sensitivity to model selection and action assumptions}
The selector sweeps preserve the direction of the routing comparison in P1, P2, P5 and P6: retained-load has lower mean cost at every tested setting of either selector threshold. For example, varying $\Delta_R$ from 0 to 0.01 changes P5's predictive-selected normalized cost from 0.575 to 0.879, versus 0.473 for retained-load. P4 favors predictive selection throughout; P3 changes ordering as the return threshold changes. The effect is therefore broader than the default cutoff, but remains process-dependent. These are diagnostic comparisons, with no test-based choice of a new threshold.

Under P2's separate case-mix shift to 65\% complex cases, the unchanged retained and predictive policies cost 0.623 and 0.699 after normalization. Intervention uncertainty matters more in P3: at actual rework reductions of 30\% and 70\%, retained-load costs 0.742 and 0.573, versus 0.671 and 0.600 for predictive selection. The ordering reverses while policies remain fixed. A response card must therefore record the effect range under which its recommendation was evaluated.

